\NativeSlot{49}{36}{%
% p022_s01

\ExamQuestion{8.}{5}%
芳芳和悠悠从圆形跑道同一地点，同时出发，相向而行，2分钟后两人相遇，芳芳每分钟跑65米，悠悠每分钟跑85米，这个圆形跑道的周长是\mbox{（\hspace{1.8em}）}米。\par
}
\NativeSlot{102}{36}{%
% p022_s02

\ExamQuestion{9.}{5}%
小华6分钟打432个字，照这样计算，他半小时打\mbox{（\hspace{1.8em}）}个。\par
}
\NativeSlot{127}{36}{%
% p022_s03

\ExamQuestion{10.}{5}%
甲、乙两地相距470千米，汽车从甲地出发行驶4小时后，还有30千米才到达终点。这辆汽车每小时行驶\mbox{（\hspace{1.8em}）}千米。\par
}
\NativeSlot{176}{76}{%
\ExamGroup{二、选择（将正确答案的序号填在括号里）。{\normalfont\kaishu\fontsize{10.8bp}{14bp}\selectfont （每题2分，共12分）}}

% p022_s04

\ExamQuestion{1.}{5}%
卡车以70千米/时的速度从甲地开往乙地，2小时后到达乙地。求甲、乙两地相距多远？列式解答时对应的数量关系式是\mbox{（\hspace{1.8em}）}。\par
\ExamOptions{A. 速度÷时间=路程}{B. 总价÷数量=单价}{C. 速度×时间=路程}{D. 单价×数量=总价}
}
\NativeSlot{303}{42}{%
% p022_s05

\ExamQuestion{2.}{5}%
一支钢笔5元钱，一个文具盒8元钱，买2支钢笔和3个文具盒，共花\mbox{（\hspace{1.8em}）}元。\par
\ExamOptions{A. 34}{B. 31}{C. 30}{D. 32}
}
\NativeSlot{378}{45}{%
% p022_s06

\ExamQuestion{3.}{5}%
北京南站到济南西站的高铁全长约406千米。如右图，G823次列车的平均速度大约是\mbox{（\hspace{1.8em}）}千米/时。\par
\ExamOptions{A. 230}{B. 203}{C. 812}{D. 820}
\par\smallskip\Needspace{4\baselineskip}\begingroup\small\setlength{\tabcolsep}{3pt}
\noindent\begin{tabularx}{\linewidth}{|*{4}{>{\centering\arraybackslash}X|}}\hline
出发站 & 出发时刻 & 到达站 & 到达时刻\\\hline
北京南站 & 12:43 & 济南西站 & 14:43\\\hline
\end{tabularx}\par
\noindent 用时2小时\par
\endgroup\smallskip
}
\NativeSlot{455}{72}{%
% p022_s07

\ExamQuestion{4.}{5}%
下列动物中，奔跑速度最快的是\mbox{（\hspace{1.8em}）}。\par
\ExamOptions{A. 羚羊每分钟约跑1440米}{B. 野驴的速度约是1200米/分钟}{C. 长颈鹿3分钟约跑3千米}{D. 东北虎1分钟约跑1300米}
}
\NativeSlot{577}{45}{%
% p022_s08

\ExamQuestion{5.}{5}%
小华和小明分别从一座桥的两端同时出发。小华的速度是70米/分，小明的速度是60米/分，经过6分钟两人相遇。这座桥长\mbox{（\hspace{1.8em}）}米。\par
\ExamOptions{A. 390}{B. 780}{C. 260}{D. 130}
}
\NativeSlot{654}{75}{%
% p022_s09

\ExamQuestion{6.}{5}%
甲乙两地相距700千米，李叔叔开车从甲地出发，去乙地送货物，每小时行驶78千米，3小时后到达服务区休息，服务区的位置可能在\mbox{（\hspace{1.8em}）}点。\par
\ExamOptions{A. ①号}{B. ②号}{C. ③号}{D. ④号}
\BookFigure{figures/p022_s09-01.tex}
}
